% Generated draft: requires mathematical and source review.
\documentclass[11pt]{article}
\usepackage[margin=25mm]{geometry}
\usepackage{fontspec}
\setmainfont{Latin Modern Roman}
\usepackage{amsmath,amssymb,mathtools,unicode-math}
\setmathfont{Latin Modern Math}
\usepackage{enumitem,needspace,xurl,hyperref}
\setlength{\parindent}{0pt}
\setlength{\parskip}{5pt}
\setlength{\emergencystretch}{3em}
\newcommand{\N}{\mathbb{N}}
\newcommand{\Z}{\mathbb{Z}}
\newcommand{\Q}{\mathbb{Q}}
\newcommand{\R}{\mathbb{R}}
\newcommand{\C}{\mathbb{C}}
\newcommand{\F}{\mathbb{F}}
\newcommand{\A}{\mathbb{A}}
\newcommand{\PP}{\mathbb{P}}
\newcommand{\E}{\mathbb{E}}
\newcommand{\eps}{\varepsilon}
\newcommand{\dd}{\,\mathrm{d}}
\begin{document}
\section*{UnsolvedMath ID: 30002372}
\textbf{Source problem code:} catalogue entry
\subsection{Definitions and mathematical statement}
Let \((M,g)\) be a polyhedral manifold with its polyhedral metric. In the usual piecewise Euclidean interpretation, the regular locus \(M_{\mathrm{reg}}\) consists of points having Euclidean metric neighborhoods, and parallel transport along loops in this locus defines a holonomy representation on a tangent space. Such a representation is irreducible if it preserves no nonzero proper real linear subspace. Write
\[
 b_2(M)=\dim_{\Q}H^2(M;\Q)
\]
for the second Betti number, using the conventional rational-coefficient interpretation. The supplied record does not specify the precise class of polyhedral metrics, the curvature convention, or the intended holonomy representation; these conventions therefore require confirmation.

\textbf{Polyhedral Frankel conjecture (supplied formulation).} For every orientable polyhedral manifold \((M,g)\) with nonnegative curvature, irreducible holonomy, and \(b_2(M)>0\), is there an integer \(n\geq 1\) and a natural integrable complex structure \(J\) on \(M\) such that
\[
 \dim_{\R}M=2n,
 \qquad (M,J)\cong_{\mathrm{bihol}}\mathbb{CP}^{n},
\]
and the given metric \(g\) is a singular K\"ahler metric with respect to \(J\)? Here \(\mathbb{CP}^{n}\) is the space of complex lines in \(\C^{n+1}\), and a biholomorphism is a holomorphic bijection with holomorphic inverse. The metric in the conclusion is the original polyhedral metric, not an independently chosen metric.

\textbf{Unresolved qualifications.} The source does not explicitly impose compactness, connectedness, completeness, or absence of boundary. None is added here. The terms ``natural holomorphic structure'' and ``singular K\"ahler metric'' are not defined in the supplied record. In the usual metric-compatible interpretation, the restriction to the regular locus would be K\"ahler: \(g\) would be Hermitian and its associated two-form \(\omega(X,Y)=g(JX,Y)\) would satisfy \(\dd\omega=0\). This regular-locus condition does not specify the required behavior at the singular set or the meaning of naturality. Consequently, the displayed assertion preserves the source target but requires clarification of these hypotheses and terms before it becomes a fully specified mathematical conjecture.
\subsection{Short English statement}
The supplied conjecture asks whether every orientable, nonnegatively curved polyhedral manifold with irreducible holonomy and positive second Betti number has a natural complex structure making it complex projective space, with its original metric singular K\"ahler. The source leaves several geometric conventions and global hypotheses unspecified. Its partially solved label and low-dimensional claims are supplied dated curation, not independently verified results.
\subsection{Sources}
\begin{itemize}
\item Local UnsolvedMath registry, permanent ID \texttt{30002372}, ``Polyhedral Frankel Conjecture,'' source problem code \texttt{catalogue entry}, collection ``Oberwolfach Reports Open Problems.'' No individual proposer is supplied; the proposed year is recorded as 2013.
\item Supplied primary provenance: \emph{Differentialgeometrie im Großen (2014)}, Oberwolfach Reports, pp.\ 1929--1974, DOI \texttt{10.4171/owr/2013/33}. \url{https://doi.org/10.4171/owr/2013/33}. The supplied background describes Oberwolfach Reports as documenting mathematical workshops and their open-problem discussions. The citation's 2014 label and its 2013 DOI component are retained without reconciliation.
\item Supplied literature source: ``Oberwolfach report on polyhedral Frankel conjecture.'' \url{https://ems.press/content/serial-article-files/46372?nt=1}.
\item Supplied literature source: ``Panov, polyhedral Kahler manifolds.'' \url{https://doi.org/10.2140/gt.2009.13.2205}.
\item Supplied literature source: ``Oberwolfach report on applications.'' \url{https://ems.press/content/serial-article-files/46468?nt=1}.
\item Literature assessment dated exactly \texttt{2026-08-21}: the record labels the problem ``partially solved'' and states that it is known in low dimensions through complex-extension results, while the general all-dimensional classification by \(\mathbb{CP}^{n}\) remains open. This is supplied dated curation derived from the OWR material, not independently verified here. The record itself requests independent verification of time-sensitive status and cited claims; it supplies neither dimension bounds nor precise hypotheses for the asserted low-dimensional results.
\item The registry labels the statement \texttt{corrected\_verified} and says that mathematical content was unchanged while notation, typography, and wording were normalized. This is supplied editorial metadata, not an independent mathematical verification of this draft.
\end{itemize}
\end{document}
